% LSPV-3 REPORT::MONOCIRCFW1::2026-10-13T11:31:50+09:00
% !TEX program = xelatex
\documentclass[11pt]{article}
\usepackage[T1]{fontenc}
\usepackage[margin=0.86in]{geometry}
\usepackage{amsmath,amssymb,amsthm,booktabs,tabularx,xurl,hyperref}
\newtheorem{proposition}{Proposition}
\newcommand{\EXACT}{\textbf{PROVED/EXACT}}
\newcommand{\KNOWN}{\textbf{KNOWN THEOREM}}
\newcommand{\COMP}{\textbf{COMPUTATIONAL}}
\newcommand{\CAND}{\textbf{CANDIDATE}}
\newcommand{\OPEN}{\textbf{OPEN}}
\newcommand{\RET}{\textbf{FALSIFIED/RETIRED}}
\newcommand{\PC}{\mathsf{PC}}
\newcommand{\PCR}{\mathsf{PCR}}
\newcommand{\IPS}{\mathsf{IPS}}
\newcommand{\Ltwo}{L_{22}}
\hypersetup{colorlinks=true,urlcolor=blue,linkcolor=blue}
\title{MONOCIRCFW1: Expanded-Monomial Hardness Does Not Lower-Bound Compact Computation}
\author{UNIVERSAL\_CORE\_REDTEAM / LSPV-3}
\date{13 October 2026}
\begin{document}
\maketitle

\begin{center}
\fbox{\parbox{0.94\linewidth}{
\textbf{Executive status: P versus NP remains OPEN.}
The all-field graph-colouring lower bound and its exact-$(2,2)$ projection
remain valid for ordinary PC/PCR with fully expanded monomial size.  This run
proves that the generic promotion from expanded support to compact arithmetic
circuit size is false.  It constructs no short circuit refutation of the
compiled graph formulas and no polynomial decider for $\Ltwo$.}}
\end{center}

\section{Exact target and claim card}
The fixed language $\Ltwo$ consists of satisfiable canonical encodings of
nonempty, proper, distinct-clause signed exact-$(2,2)$ 3-CNFs; malformed or
out-of-domain strings are rejected.  The desired direct theorem is
\[
 \forall A\ \forall k\ \forall C>0\ \forall N_0\ \exists G\quad
 [G\in\operatorname{Dom}(L_{22})\land |G|\ge N_0\land
   T_A(G)>C|G|^k],                                      \tag{1}
\]
where $A$ ranges over total deterministic uniform deciders of this one
language.  Inputs may depend on $A$.

The falsification target in this run was the weaker-looking transfer
\[
 \text{exponential expanded-monomial PC/PCR size}
 \Longrightarrow
 \text{superpolynomial compact algebraic computation or decision time}.
                                                               \tag{2}
\]
The source realization is the fully expanded PC/PCR proof.  The proposed
target realizations are arithmetic circuits, circuit-based algebraic proofs,
or a one-bit machine execution.  They are different objects; no
\textsc{same\_realization} theorem was supplied.

\section{Predecessor and corpus review}
\textbf{CORPUS REVIEWED.}
The immediately preceding LSPV-2 artifact
\path{research/agents/lspv-2/2026-10-13_1122_LSPV-2_COMPOSEPROJ1.tex}
was read back from GitHub with blob
\texttt{5b62202ab07c5680a1f05f4f8584f91993126e1b}; it is substantive and
has both required sentinels.  No predecessor recovery was needed.

This run surveyed the current GitHub tree; read root \path{main.tex},
\path{AGENTS.md}, \path{CODEX.md}, the full pnp-research skill and red-team
protocol, all current task specifications surfaced under \path{tasks/},
actual \path{main4.tex} and \path{main5.tex}, the newest verified MAIN
\path{research/main/2026-10-13_0948_MAIN_BLACKBOXQ1.tex}, and materially
relevant LSPV-1--4 reports.  It read the HEC/VSC-HEC/G-PPC-HEC/CRS-PPC/CCSPC
guards surfaced by repository search, the SINGLE22 and EARLY22 governance
notes, the exact Drive recovery snapshot, the graph compiler source
\path{research/agents/lspv-4/2026-10-13_0035_LSPV-4_PC22_ALLFIELD1.tex}
(blob \texttt{4166635f614e8baa6a94bff5cf47e76815e971f5}),
\path{experiments/review/universal22.py},
\path{experiments/review/audit_universal22.py}, COMPOSEPROJ1's code/result,
TRACEFRONT1, TRACESIMQ1, IMAGE22COMP1, and EARLY22RED1.

\textbf{DUPLICATION / CONTRADICTION AUDIT.}
The graph-compiler report already warns that expanded monomial count and
arithmetic-circuit size differ, using the same product family below.  This
run does not claim that warning as new.  Its new contribution is the formal
transfer retirement, primary-literature cross-check, hostile-control program,
and root reconciliation.  A genuine stale statement was found in the LSPV-3
root region: it limited the continuation ledger to source-clause punctures.
The later verified TRACEFRONT1 ledger records that ``21 P C B'' covers every
one-clause target puncture and target restriction by bounded local work plus
affine continuation.  The subsidiary statement is corrected; the compiler,
minimality, and Resolution projection are unaffected.

\section{Primary-literature and barrier preflight}
The two active obstacles are (i) a representation mismatch between sparse
expanded PC/PCR lines and compact arithmetic circuits, and (ii) the absence
of a theorem converting an arbitrary one-bit $L_{22}$ computation into a
hard algebraic proof in the same realization.

Mik\v{s}a--Nordstr\"om define PC/PCR proof size as the number of monomials
after every polynomial is expanded; their degree-to-size machinery is in that
model.  Conneryd--de Rezende--Nordstr\"om--Pang--Risse prove linear PC degree,
hence strongly exponential PC size, for sparse random graph 3-colouring over
every field.  This is the exact proof-system endpoint used by PC22\_ALLFIELD1.

Grochow--Pitassi introduce the circuit-based Ideal Proof System and explicitly
connect its lower bounds to algebraic circuit lower bounds.  Forbes--Shpilka--
Tzameret--Wigderson study restricted IPS circuit classes and emphasize that
degree or sparsity lower bounds are not automatically circuit-complexity lower
bounds.  Thus circuit representation is not a harmless re-encoding of the
existing PC/PCR measure; it is a stronger model requiring a new theorem.

\begin{center}
\begin{tabularx}{0.96\linewidth}{@{}p{0.23\linewidth}X@{}}
\toprule
Classification & Disposition\\
\midrule
THEOREM TRANSFER & The graph lower bound transfers through the audited
source/properization/compiler substitution to PC/PCR degree and fully
expanded monomial size only.\\
ANALOGY / CANDIDATE & Circuit-based IPS suggests the right stronger algebraic
model, but a lower bound for the compiled graph family in that model is not
known here.\\
FAILED TRANSFER & Expanded PC/PCR monomial size, VLSI routing, communication
cuts, Schur-square statistics, or cryptographic binding do not lower-bound an
arbitrary uniform one-bit decider without a model-matched theorem.\\
\bottomrule
\end{tabularx}
\end{center}

\section{The compact-representation firewall}
\begin{proposition}[Exact counterexample to generic support transfer]
For every field $\mathbb F$ and $t\ge1$, the multilinear polynomial
\[
                  P_t(x)=\prod_{i=1}^t(1+x_i)                 \tag{3}
\]
has exactly $2^t$ distinct monomials in expanded form, yet has a division-free
arithmetic circuit with $t$ addition gates and $t-1$ multiplication gates.
It can be balanced to depth at most $1+\lceil\log_2t\rceil$.
Consequently, expanded monomial count does not lower-bound arithmetic-circuit
size even up to any fixed polynomial.
\end{proposition}
\begin{proof}
Expansion chooses independently either $1$ or $x_i$ from every factor.  Each
subset of $[t]$ produces a different squarefree monomial, with coefficient
one in every field.  Computing each factor costs one addition and multiplying
the factors in a binary tree costs $t-1$ multiplications.  The ratio
$2^t/(2t-1)$ is superpolynomial.
\end{proof}

The separation is robust to the requested wrappers.  Variable permutations
leave both complexities unchanged.  Complementing any input is implemented
by one affine gate per changed variable.  Over characteristic two,
complementing every variable transforms (3) into
\(
  \prod_i(1+(1-x_i))=\prod_i x_i,
\)
whose expanded support has size one.  Independent products multiply expanded
supports while their circuits compose with one multiplication gate.  Padding,
COPY, INTERP, and RECOMP can preserve or shrink the compact representation
without preserving an expanded support count.

Therefore \texttt{EXPANDED-MONOMIAL-TO-COMPACT-COMPUTATION1} is
\RET.  This does \emph{not} show that the compiled graph-colouring formulas
have short IPS or arithmetic-circuit refutations.  It shows that the PC/PCR
lower bound alone cannot rule them out.  The smallest repair is to prove a
lower bound in the declared circuit-represented proof model itself.  To infer
(1) from such a result would still require either a direct machine lower bound
or a separately proved charged machine-to-proof theorem; the retired universal
trace/extraction route cannot be assumed.

\section{Bounded reproducible audit}
The independent script
\path{experiments/review/expanded_monomial_circuit_firewall.py}
(GitHub commit \texttt{6100ac7df7296acba0cffad846c1269f52d9029e}, blob
\texttt{d130587a5ed8a9eede52a5ba864cafbbc8435172}, SHA-256
\texttt{040d33951d0e2563b22d8920143a13b48367004362a30703e74df3a04a95f850})
produced
\path{results/raw/review/2026-10-13_1131_expanded_monomial_circuit_firewall.json}
(GitHub commit \texttt{b62d0ef736590afc91516bb190923f6b9e0769b8}, blob
\texttt{4fbc5eeb2cd1d9e7cbe2adb04057b0d4c382145d}).

One 12.01-second invocation made 302 checks with no failure.  It expanded
(3) exhaustively only for $1\le t\le18$, checked compact and expanded
evaluation over fields $\mathbb F_p$ for $p=2,3,5,7,11$ under identity,
alternating phase, and all-phase substitutions, and checked independent
products.  At $t=18$ it observed 262,144 expanded monomials versus 35
arithmetic operation gates; the all-complemented characteristic-two support
was one.  The semantic digest is
\texttt{19386eec6f15c462e029600ebee54f5e5d69a859c1638a8671b89b98aea965eb}.

The same run exercised actual compact one-bit controls at size 64: SCC-based
2SAT, Horn forward chaining, Gaussian elimination of rank 63 with consistent
and inconsistent cases, augmenting-path bipartite matching of size 64, and
width-one path dynamic programming.  These controls are not algorithms for
the graph-colouring source or $L_{22}$; they prevent a representation metric
from falsely charging known compact computations.  No $2^{144}$ assignment
enumeration and no asymptotic inference from finite timing were made.

\section{Barrier and proof-status audit}
\textbf{Relativization.}  The firewall is algebraic and survives semantic
oracle replacement; an oracle can only widen the mismatch.  It provides no
nonrelativizing hinge for (1).

\textbf{Algebrization.}  Compact algebraic access is exactly the hostile
replacement that breaks the generic transfer.  The all-field PC/PCR theorem
does not by itself cross this barrier.

\textbf{Natural Proofs.}  No broad constructive Boolean circuit property is
proved.  If the repair attempts a general circuit lower bound, Natural-Proofs
hygiene becomes mandatory.  The fixed-language goal does not require an
unnecessary general circuit lower bound.

\textbf{Arbitrary exponents and nonuniformity.}  A polynomial $L_{22}$ decider
may have any fixed exponent depending on the machine.  Nothing in a
$2^{\Omega(v)}$ expanded-PC lower bound dominates that exponent in the
machine model without a transfer.  Nonuniform compact circuits are an even
stronger hostile representation and cannot be excluded by the current proof.

\textbf{One-bit output and proof search.}  Verifying or lower-bounding an
explicit refutation is different from finding one bit.  An arbitrary decider
need not emit, store, or expose the PC/PCR lines.  TRACEFRONT1 additionally
gives a total correct front-door decider fast on a recognizable
Resolution-hard sequence, falsifying a universal trace-to-Resolution
promotion.  The same quantifier error must not be reintroduced for PC/PCR.

\section{EARLY22 and SINGLE22 acknowledgments}
\paragraph{EARLY22-INTEGRATION1.}
Sources actually read were the governance memo, the exact Drive snapshot
(SHA-256 \texttt{b6be320498d5e848ba4d6a023bf03b18d8085ebabec1d92368180a9739025dac}),
the original LBO3SAT-0031/0032 and PROOF-0034 excerpts, EARLY22RED1, and the
later TRACEFRONT1 supersession ledger.  Disposition: \textbf{INTEGRATED WITH
SUBSIDIARY SCOPE REPAIR}.  The exact manuscript locations are Sections 2 and
7 here and the reconciled LSPV-3 root region.  The compiled minimally-UNSAT
exact-$(2,2)$ Resolution family, monotone cubic-pair map, all-field graph PC
family, and continuation collapse remain separate entries.  Next check:
read any proposed compact-circuit proof lower bound against the exact graph
encoding before it is used; do not reopen one-clause continuation.

\paragraph{SINGLE22-SEPARATION1.}
The fixed-language memo, PC22 graph-colouring report at its full path and
blob, both compiler scripts, COMPOSEPROJ1, and IMAGE22COMP1 were read.
Disposition: \textbf{INTEGRATED}.  The exact arrow advanced is negative:
the representation-only route from all-field expanded PC/PCR size to compact
computation is retired.  The graph adapter and proof projection remain
positive inputs; no family broadening is demanded.  The next proof obligation
is either a direct whole-machine lower bound for (1), or a genuinely new
model-matched compact-proof lower bound plus a noncircular charged implication
from the relevant computation.  Neither is supplied here.

\section{Required reconciliation fields}
\textbf{WHAT CHANGED RELATIVE TO ROOT MAIN.}
The LSPV-3 region now records the all-one-clause continuation supersession and
the compact-representation firewall.  Executive synthesis outside the owned
region is unchanged.

\textbf{ADOPTED / REJECTED / RETIRED PRIOR COMPONENTS.}
Adopted: $L_{22}$ NP-completeness, graph compiler/projection, all-field PC/PCR
degree and expanded-size theorem, exact-(2,2) compiler/minimality/Resolution
projection, and 21 P C B's scoped continuation collapse.  Rejected as
unproved: a compact IPS lower bound for this family.  Retired: generic
expanded-monomial-to-circuit/runtime transfer and universal trace-to-proof
promotion.

\textbf{STALE OR OVERSTATED CLAIMS FOUND.}
The LSPV-3 root restriction of continuation to source punctures and its
proposed non-source-puncture next check were stale after 21 P C B.  They are
replaced by the verified bounded-local-plus-affine scope.  No positive
compiler theorem is retired.

\begin{thebibliography}{9}
\bibitem{conneryd} J. Conneryd, S. F. de Rezende, J. Nordstr\"om, S. Pang,
and K. Risse, ``Graph Colouring Is Hard on Average for Polynomial Calculus
and Nullstellensatz,'' arXiv:2503.17022, 2025.
\bibitem{mn} M. Mik\v{s}a and J. Nordstr\"om, ``A Generalized Method for
Proving Polynomial Calculus Degree Lower Bounds,'' CCC 2015,
doi:10.4230/LIPIcs.CCC.2015.467.
\bibitem{gp} J. A. Grochow and T. Pitassi, ``Circuit Complexity, Proof
Complexity, and Polynomial Identity Testing,'' arXiv:1404.3820; JACM 2018.
\bibitem{fstw} M. A. Forbes, A. Shpilka, I. Tzameret, and A. Wigderson,
``Proof Complexity Lower Bounds from Algebraic Circuit Complexity,''
arXiv:1606.05050, CCC 2016.
\end{thebibliography}

\end{document}
% END_OF_VERIFIED_REPORT::LSPV-3::MONOCIRCFW1
